\color{unidarkgreen}{\subsection{Cигнализация}}
\color{black}

\begin{enumerate}

\item
Устройство формирует выходное воздействие на сигнализацию «Срабатывание» по схеме «ИЛИ» в следующих случаях:
\begin{itemize}
\item Срабатывание резервных защит;
\item Срабатывание ТЗО;
\item Срабатывание внешних защит;
\item Срабатывание АУ;
\item Срабатывание ЛД;
\item Срабатывание ЗНР;
\item Срабатывание ЗНФ В1 и ЗНФ В2.
\end{itemize}

\item
Устройство формирует выходное воздействие на сигнализацию «Неисправность» по схеме «ИЛИ» в следующих случаях:
\begin{itemize}
\item Неисправность цепей напряжения;
\item Зависание сигнала внешнего отключения;
\item Неисправность силового выключателя;
\item Нарушение режима по испытательным блокам и/или переключателям в выходных цепях;
\item Потеря GOOSE-сообщения;
\item Срабатывание внешней функции на сигнализацию.
\end{itemize}

\item
Алгоритм работы логики логики предупредительной сигнализации приведен на рисунке \ref{pic:SIGN}.

\medskip
\begin{figure}[H]
\centering
\includegraphics[width=1\textwidth,height=1\textheight,keepaspectratio]{\fbpath/07.02 РЗАТ/91 Сигнализация/img/SIGN.pdf}
\captionof{figure}{Функциональная схема логики блока сигнализации}\label{pic:SIGN}
\end{figure}

\end{enumerate}